\section{The cancelled openness remark in cote 34}
\label{sec:q03}

The literal assertion in cote 34, scan 064, Remark 1.17, that
\(\overline{[\pi_0,H]}\) is open in \(\pi'\), is false.
There are two separate issues: the commutators lie in the geometric
kernel \(H\), and a quotient \(\pi_0\) acts on a nonabelian kernel only
through outer automorphisms unless an action or a section is specified.
We prove the natural precise corrections, including the soluble
nonabelian statement.  All openness conclusions below concern \(H\).
The exact intended wording of the crossed-out sentence remains a
question about the original image; a forced mathematical correction is
not a verification of handwriting.

\subsection{The theorem already available in the source}

Let \(k\) be finitely generated over its prime field, let \(X/k\) be a
normal geometrically connected scheme locally of finite type, and let
\(\ell\ne\operatorname{char}(k)\).  For an open subgroup
\(\pi\subseteq\pi_1(X)\), write
\[
1\longrightarrow I\longrightarrow\pi\longrightarrow\Gamma
\longrightarrow1,
\qquad
I=\pi\cap\pi_1(X_{k^{\mathrm{sep}}}).
\]
Thus \(\Gamma\), denoted \(\pi_0\) in the source, is the arithmetic
image.  For a profinite group \(J\), let \(J(\ell)\) denote its maximal
\emph{abelian} pro-\(\ell\) quotient.  Commutator subgroups and
coinvariant relations are always closed topologically.

The background theorem used in this section is precisely cote 34,
scan 053, Theorem 1.1:
\[
 I(\ell)\text{ is a finitely generated }\mathbf Z_\ell\text{-module},
 \qquad I(\ell)_\Gamma\text{ is finite}.
 \tag{Q3-B01}
\]
Its quantifier over \emph{every open subgroup} \(\pi\) is essential.
We carry this established theorem as background, rather than claiming
to have reconstructed the omitted proof pages 065--073.  The supplied
scans 053--063 were read before Remark 1.17; scan 074 was also checked.

Let \(N_0\subseteq I\) be closed and normal in \(\pi\), and suppose
\[
H=I/N_0
\]
is pro-\(\ell\).  Put \(P=\pi/N_0\).  The extension under discussion is
\[
1\longrightarrow H\longrightarrow P\longrightarrow\Gamma
\longrightarrow1. \tag{Q3-E}
\]
Here \(P\) is the \(\pi'\) of the cancelled remark.  Distinguishing this
symbol from the abelian quotient extension used earlier in Theorem 1.1
avoids an ambiguity in the source.

\begin{lemma}[Q3-03: relative finite abelianization]
\label{q3:relative}
For every open subgroup \(P_1\subseteq P\), let \(J=P_1\cap H\).
Then the coinvariant group
\[
(J^{\mathrm{ab}})_{P_1}
\]
for the conjugation action is finite.
\end{lemma}

\begin{proof}
Let \(\pi_1\) be the inverse image of \(P_1\) in \(\pi\).  It is open
in \(\pi_1(X)\).  Its geometric kernel
\(I_1=\pi_1\cap I\) maps onto \(J\): a lift in \(\pi_1\) of an
element of \(J\) has trivial arithmetic image and so belongs to
\(I_1\).  Since \(J\) is pro-\(\ell\), the composite
\[
I_1\twoheadrightarrow J\twoheadrightarrow J^{\mathrm{ab}}
\]
factors through \(I_1(\ell)\).  It is equivariant for \(\pi_1\).
On \(I_1(\ell)\) the conjugation action factors through the arithmetic
image \(\Gamma_1\) of \(\pi_1\).  Consequently there is a surjection
\[
I_1(\ell)_{\Gamma_1}\twoheadrightarrow
(J^{\mathrm{ab}})_{P_1}.
\]
The source is finite by Theorem 1.1 applied to the open subgroup
\(\pi_1\).  This proves the assertion.  No section or solvability
assumption was used.
\end{proof}

\subsection{Literal refutation and the canonical correction}

\begin{proposition}[Q3-01: the stated ambient group is impossible]
\label{q3:literal}
Even when a continuous section of \((\mathrm{Q3\text{-}E})\) is given
and \(H\) is soluble of finite derived length,
\(\overline{[\Gamma,H]}\) need not be open in \(P\).
\end{proposition}

\begin{proof}
Take \(k=\mathbf F_q\), \(X=\operatorname{Spec}k\),
\(\pi=P=\Gamma=G_k=\widehat{\mathbf Z}\), and \(I=H=1\).
The trivial geometric quotient is a soluble pro-\(\ell\) group.
Its commutator subgroup with \(\Gamma\) is trivial, hence is not
open in the infinite profinite group \(P\).  These data meet every
displayed hypothesis.

There is also a nontrivial geometric example.  Take
\(X=\mathbf G_m/\mathbf F_2\), \(\ell=3\), and the rational point
\(1\).  The Kummer tower gives
\[
H=\mathbf Z_3(1),
\]
and the rational point gives a section.  Arithmetic Frobenius acts
on \(H\) by multiplication by \(2\).  Therefore its commutators
already contain \((2-1)H=H\), so
\(\overline{[\Gamma,H]}=H\).  But
\(P/H=\widehat{\mathbf Z}\) is infinite, and \(H\) is not open in
\(P\).
\end{proof}

For a general extension \((\mathrm{Q3\text{-}E})\), conjugation gives
a canonical homomorphism
\[
\Gamma\longrightarrow\operatorname{Out}(H).
\]
Choosing different lifts in \(P\) changes the action on \(H\) by
inner automorphisms.  Thus \([\Gamma,H]\), construed as the
commutator subgroup for an actual action, needs additional data.
The subgroup \([P,H]\) does not.

\begin{proposition}[Q3-04: canonical openness]
\label{q3:canonical}
For every pro-\(\ell\) quotient \(H\) in
\((\mathrm{Q3\text{-}E})\), the subgroup
\[
C=\overline{[P,H]}
\]
is open in \(H\).  No solvability assumption is needed.
\end{proposition}

\begin{proof}
Since \(H\) is normal in \(P\), the closed subgroup \(C\) is normal
in \(P\).  It contains \(\overline{[H,H]}\), so \(H/C\) is abelian.
By definition, \(P\) acts trivially on this quotient.  Hence the
quotient map \(I\to H/C\) factors through
\[
I(\ell)_\Gamma.
\]
This last group is finite by Theorem 1.1, and therefore \(H/C\) is
finite.  A closed subgroup of finite index in a profinite group is
open.
\end{proof}

Equivalently,
\[
H/\overline{[P,H]}=(H^{\mathrm{ab}})_\Gamma.
\]
This correction is exactly the quotient already controlled by
Theorem 1.1.  It does not require a new theorem about purity in
arbitrary dimension.

\subsection{The action-dependent soluble statement}

Suppose now that a continuous section \(s:\Gamma\to P\) is supplied.
Write
\[
\alpha_\gamma(h)=s(\gamma)h s(\gamma)^{-1},\qquad
\delta_\gamma(h)=h^{-1}\alpha_\gamma(h),
\]
and define
\[
N=\overline{\left\langle
 \delta_\gamma(h):h\in H,\ \gamma\in\Gamma
\right\rangle}.
\]
The word soluble in the next result means that the \emph{closed}
derived series terminates after finitely many steps.

\begin{proposition}[Q3-05: normality and the FAb quotient]
\label{q3:fab}
The subgroup \(N\) is normal in \(P\).  For \(Q=H/N\), every open
subgroup \(V\subseteq Q\) has finite abelianization.
\end{proposition}

\begin{proof}
The normality assertion must be checked before forming the
nonabelian quotient.  For \(h,k\in H\), direct multiplication gives
\[
\delta_\gamma(hk)
=k^{-1}\delta_\gamma(h)k\,\delta_\gamma(k).
\]
Thus
\[
k^{-1}\delta_\gamma(h)k
=\delta_\gamma(hk)\delta_\gamma(k)^{-1}.
\]
The subgroup generated by the \(\delta_\gamma(h)\), and therefore
its closure \(N\), is normal in \(H\).  For \(\eta\in\Gamma\),
\[
\alpha_\eta\bigl(\delta_\gamma(h)\bigr)
=\delta_{\eta\gamma\eta^{-1}}\bigl(\alpha_\eta(h)\bigr).
\]
It follows that \(N\) is stable under \(s(\Gamma)\), and hence is
normal in \(P=H\rtimes s(\Gamma)\).

The action of \(\Gamma\) on \(Q=H/N\) is trivial.  The given section
therefore identifies the quotient extension with the direct product
\[
P/N\simeq Q\times\Gamma.
\]
Let \(V\) be an arbitrary open subgroup of \(Q\), not necessarily
normal, and let \(P_V\) be the inverse image of \(V\times\Gamma\)
in \(P\).  The subgroup \(P_V\) is open, and
\[
J_V=P_V\cap H\twoheadrightarrow V.
\]
Conjugation by \(V\times\Gamma\) acts trivially on
\(V^{\mathrm{ab}}\): the first factor acts by inner automorphisms,
and the second acts trivially already on \(V\).  Therefore
\(V^{\mathrm{ab}}\) is a quotient of
\[
(J_V^{\mathrm{ab}})_{P_V}.
\]
Lemma~\ref{q3:relative} makes this group finite.  Thus every open
subgroup \(V\) of \(Q\) has finite abelianization, as asserted.
\end{proof}

A profinite group with the conclusion of
Proposition~\ref{q3:fab} is called \emph{FAb}.  Finiteness of
\(Q^{\mathrm{ab}}\) alone would be insufficient for the next
argument; the quantifier over all open subgroups is what permits
iteration.

\begin{lemma}[Q3-06: soluble FAb groups]
\label{q3:soluble}
A profinite FAb group of finite derived length is finite.
Consequently, if \(H\) in Proposition~\ref{q3:fab} has finite
derived length, then \(N=\overline{[s(\Gamma),H]}\) is open in \(H\).
\end{lemma}

\begin{proof}
For an FAb group \(Q\), let
\[
D^0Q=Q,\qquad
D^{r+1}Q=\overline{[D^rQ,D^rQ]}.
\]
The quotient \(Q^{\mathrm{ab}}\) is finite, so \(D^1Q\) is open.
If \(D^rQ\) is open in \(Q\), FAb says that its abelianization is
finite.  Hence \(D^{r+1}Q\) is open in \(D^rQ\), and therefore in
\(Q\).  Induction makes every \(D^rQ\) open.  Finite derived length
gives \(D^mQ=1\) for some finite \(m\).  Thus \(Q\) is compact and
discrete, and hence finite.

In Proposition~\ref{q3:fab}, the quotient \(Q=H/N\) has no greater
derived length than \(H\).  It is FAb by that proposition, so it is
finite.  This is exactly the openness of \(N\) in \(H\).
\end{proof}

\begin{remark}
Replacing finite derived length by prosolvability would erase a
necessary hypothesis.  Every pro-\(\ell\) group is already
prosolvable, whereas the proof above requires a finite terminal
derived index.  No such extension is asserted here.
\end{remark}

\subsection{An intrinsic variant without a section}

\begin{proposition}[Q3-07: trivial outer arithmetic action]
\label{q3:outer}
Suppose \(H\) in \((\mathrm{Q3\text{-}E})\) has finite derived
length and the canonical map
\(\Gamma\to\operatorname{Out}(H)\) is trivial.  Then \(H\) is finite.
\end{proposition}

\begin{proof}
Let \(C=C_P(H)\), the closed centralizer of \(H\) in \(P\).
Triviality of the outer action means that every element of \(P\)
acts on \(H\) by conjugation by an element of \(H\).  Thus
\[
P=HC.
\]
Let \(D_i\) be the closed derived series of \(H\).  We prove
inductively that \(D_i\) is open in \(H\).  The claim holds for
\(D_0=H\).  Suppose it holds for \(D_i\).  The group \(D_i\) is
characteristic in \(H\), hence normal in \(P\).  The product
\[
R_i=D_i C
\]
is a closed subgroup of \(P\).  Since \(P=HC\),
\[
[P:R_i]\le[H:D_i]<\infty,
\]
so \(R_i\) is open.

Choose an open normal subgroup \(U_i\subseteq P\) such that
\(U_i\cap H\subseteq D_i\).  Such a subgroup exists because \(D_i\)
is an open neighbourhood of \(1\) in the subspace \(H\), and open
normal subgroups form a neighbourhood basis at \(1\) in \(P\).
Since \(D_i\) is normal in \(P\), \(D_iU_i\) is open and
\[
(D_iU_i)\cap H=D_i.
\]
Put
\[
P_i=R_i\cap D_iU_i.
\]
It is open in \(P\), contains \(D_i\), and satisfies
\(P_i\cap H=D_i\).  Every element of \(P_i\subseteq D_iC\) acts on
\(D_i\) by conjugation by an element of \(D_i\), because \(C\)
centralizes \(H\).  Its action on \(D_i^{\mathrm{ab}}\) is
therefore trivial.  Lemma~\ref{q3:relative}, applied to \(P_i\),
shows that \(D_i^{\mathrm{ab}}\) is finite.  Hence \(D_{i+1}\)
is open in \(D_i\), and therefore open in \(H\).

The induction reaches \(D_m=1\) after finitely many steps.  Thus
\(1\) is open in \(H\), proving that \(H\) is finite.
\end{proof}

This is an optional intrinsic result about the extension; it does
not identify the intended wording of Remark 1.17.

\subsection{The exact weight input and the role of Weil I}

\begin{lemma}[Q3-02: a Frobenius without eigenvalue \(1\)]
\label{q3:weights}
Let \(M\) be a finitely generated \(\mathbf Z_\ell\)-module with
continuous action by a profinite group \(G\).  If some \(g\in G\)
has no eigenvalue \(1\) on \(M\otimes_{\mathbf Z_\ell}\mathbf Q_\ell\),
then \(M_G\) is finite.
\end{lemma}

\begin{proof}
The module \(M/(g-1)M\) is finitely generated over
\(\mathbf Z_\ell\).  Its tensor product with \(\mathbf Q_\ell\)
vanishes because \(g-1\) is invertible after tensoring.  Thus it is
a finitely generated torsion \(\mathbf Z_\ell\)-module and is
finite.  It surjects onto \(M_G\), so \(M_G\) is finite.
\end{proof}

If \(g\) is a finite-field Frobenius with pure eigenvalues of a
nonzero weight, the eigenvalue \(1\) is excluded.  Replacing
geometric Frobenius by arithmetic Frobenius inverts eigenvalues
and reverses weights, without affecting that exclusion.
Semisimplicity is not needed.

For a smooth projective curve \(C/\mathbf F_q\), Weil I,
\cite[(1.6)--(1.7)]{WeilI}, gives weight \(1\) on
\(H^1(C_{\overline{\mathbf F}_q},\mathbf Q_\ell)\), hence weight
\(-1\) on its dual.  The identification
\[
H^1(C_{\overline{\mathbf F}_q},\mathbf Q_\ell)
=\operatorname{Hom}_{\mathrm{cont}}
 \bigl(\pi_1(C_{\overline{\mathbf F}_q}),\mathbf Q_\ell\bigr)
\]
follows from the classification of finite constant-group torsors
and passage to the inverse limit.  This supplies the
no-eigenvalue-\(1\) input for Lemma~\ref{q3:weights}.

The spreading and specialization mechanism is explicit in cote
34, scan 074, formulas (2.5.4)--(2.5.6): extend the smooth
projective curve over a finitely generated normal ring, compare
the lisse \(H^1\)-fibres, and specialize at a closed point with
finite residue field.  One Frobenius without eigenvalue \(1\)
suffices.  Scan 054 announces the reduction to curves, and scan
053 states the resulting Theorem 1.1.  This is why the curve
case contained in Weil I is sufficient for the established input
used here.  We are not asserting that smooth projective purity
alone, without those geometric reductions, proves an arbitrary
normal-scheme fundamental-group statement.

An independently accessible published account of the same
mechanism is \cite{KatzLang}: Theorem 1 and its proof in
Section II, followed by Theorems \(1(\mathrm{bis})\) and
\(1(\mathrm{ter})\).  The proof passes through Artin good
neighbourhoods and generalized Jacobians to abelian varieties
and tori over finite fields; its exact cited locators for the
neighbourhood reduction are SGA 4 XI 3.1 and 3.3.  This confirms
the curve/Jacobian/spreading mechanism, but is not substituted
without further argument for the full normal-scheme formulation
of the supplied Theorem 1.1.

For comparison, \cite[(3.3.5), p.~206]{WeilII} supplies a broader
direct weight estimate: on a smooth variety over a finite field,
\(H^1\) has weights at least \(1\), so its dual has weights at
most \(-1\).  Lemma~\ref{q3:weights} then applies.  This Weil II
route is unnecessary for
Propositions~\ref{q3:canonical}, \ref{q3:fab}, and
\ref{q3:outer}; it records accurately the distinction between
smooth projective purity and a nonproper weight estimate.

The dependency chain for the corrected statements is therefore
\[
\begin{gathered}
\text{curve RH, contained in Weil I}
+\text{ the geometric reductions used in Theorem 1.1}\\
\Longrightarrow
\text{Theorem 1.1 for every open }\pi
\Longrightarrow \text{Q3-03}\\
\Longrightarrow
\text{Q3-04, or Q3-05--Q3-06, or Q3-07}.
\end{gathered}
\]
None of these implications uses a higher-dimensional purity
theorem, a Hodge standard conjecture, or the assumed arithmetic
RH/Dirichlet-GRH input.  The literal assertion in
Proposition~\ref{q3:literal} remains false under all such
additional hypotheses.

\subsection{The locally finite type qualification}

\begin{lemma}[Q3-08: reduction to a nonempty affine open]
\label{q3:affine}
Under the exact-sequence framework stated in Theorem 1.1, its
finite-generation and finite-coinvariant conclusions for normal
geometrically connected finite-type schemes imply the same
conclusions for normal geometrically connected schemes locally
of finite type.
\end{lemma}

\begin{proof}
A connected normal locally noetherian scheme is integral:
normality makes distinct irreducible components disjoint, and
local noetherianity makes each component locally open.
Connectedness leaves a single component.  Choose a nonempty
affine open \(U\subseteq X\).  It is finite type over \(k\).

If \(X'\to X\) is a connected finite étale cover, then \(X'\)
is normal and connected, hence integral.  Its inverse image of
\(U\) is a nonempty open subset of an integral scheme and is
therefore connected.  The connected-cover criterion
\cite[V, Proposition 6.9]{SGA1} gives a surjection
\[
\pi_1(U)\twoheadrightarrow\pi_1(X).
\]
Normality is preserved by separable algebraic field extension.
Applying the same argument over \(k^{\mathrm{sep}}\), where the
assumed geometric connectedness applies, gives
\[
\pi_1(U_{k^{\mathrm{sep}}})
\twoheadrightarrow\pi_1(X_{k^{\mathrm{sep}}}).
\]

For an open subgroup \(\pi\subseteq\pi_1(X)\), let \(\pi_U\)
be its inverse image in \(\pi_1(U)\).  Its geometric kernel
\(I_U\) maps onto \(I\), and both arithmetic images are the same
group \(\Gamma\).  Consequently
\[
I_U(\ell)\twoheadrightarrow I(\ell),\qquad
I_U(\ell)_\Gamma\twoheadrightarrow I(\ell)_\Gamma.
\]
The finite-type conclusion for \(U\) therefore transfers finite
generation and finiteness of coinvariants to \(X\).
\end{proof}

The published exact-sequence locator is
\cite[IX, Theorem 6.1]{SGA1}, originally p.~253
(p.~195 in the re-edition); that statement includes a
quasi-compactness hypothesis.  The reference to Exposé VIII in
scan 053 does not match this published locator and should be
flagged as a citation candidate, possibly reflecting an earlier
numbering.  Lemma~\ref{q3:affine} does not silently apply the
quasi-compact theorem to a non-quasi-compact scheme: it uses the
exact-sequence framework already supplied in Theorem 1.1 and
proves the surjections needed to transfer its conclusions.
The underlying Galois-category equivalence is
\cite[V, Theorem 4.1]{SGA1}.

\subsection{Disposition of the question}

The literal row is refuted by Q3-01.  The canonical correction
\(\overline{[P,H]}\) open in \(H\) is Q3-04.  With a specified
continuous section and finite derived length,
\(\overline{[s(\Gamma),H]}\) open in \(H\) is Q3-05--Q3-06.
The intrinsic trivial-outer-action variant is Q3-07.
These statements are settled within the established
fundamental-group and curve-RH framework.  The remaining textual
question is which correction, if any, the author intended in the
cancelled Remark 1.17.

